\section{Failure conditions}
\label{sec:gauge-failure-conditions}

The finite chain is exact on its stated premises, and each step has a
registered negative control. The Lie-type theorem fails if a faithful
compact commutator-closed twelve-port response with inner carrier action is
exhibited outside
\(\mathfrak u(1)\oplus\mathfrak{su}(2)\oplus\mathfrak{su}(3)\); the
conditional matter results fail if the incidence identity, the target-blind
response, the equivariant lift, the anomaly scan, the tensor-character
enumeration or the refinement persistence receipt fails its control. The
categorical route applies only inside its combined strictification,
trivial-holonomy, localization and refinement and fibre receipts.

On the executable model the branch is tested by computing the complete
accessible current of a specified law, with every admitted control and
apparatus service in the menu. A complete accessible response that is not
faithful, closed and twelve-dimensional with inner carrier action excludes
that law's realization of A1 and A2, as the seventeen-dimensional exchange
response excludes its grammar. A physical attachment that forced the carrier
rechartings to be spatial rotations of an interacting relativistic field while
they remain internal would be constrained by the Coleman--Mandula theorem.
The particle paper~\cite{ophparticles} tests whether localized defects of the
same transport carry holonomy valued in this structure.
